% -----------------------------------------------------------------------------
% One-loop QED: pole, charge and vertex renormalization as one physical chain
% -----------------------------------------------------------------------------

树级 QED 只有一个电子--光子局域顶角，但一圈以后出现三个彼此不同的一粒子不可约基本核：电子自能 \(\Sigma\)、光子真空极化 \(\Pi^{\mu\nu}\) 与电磁顶角修正 \(\Lambda^\mu\)。它们分别对应三类可测信息：电子传播子的单粒子极点确定物理质量与外腿归一化；长波长电磁交换确定低能电荷；在壳电磁顶角的额外张量结构确定磁矩等自旋响应。紫外发散出现在裸参数与这些可测定义之间的微扰关系中，质量、电荷和场强反项由相应的极点、留数与低能归一化条件固定。

对以特征物理动量 \(Q\) 探测、且没有大对数重求和的过程，固定一圈截断的两个控制量为
\begin{equation*}
\epsilon_{\rm 1L}\equiv\frac{\alpha(\mu)}{\pi},
\qquad
\epsilon_{\log}\equiv
\frac{\alpha(\mu)}{\pi}
\left|\ln\frac{Q^2}{m_{\mathrm e}^2c^2}\right|.
\end{equation*}
当 \(\epsilon_{\rm 1L}\ll1\) 且 \(\epsilon_{\log}\ll1\) 时，一圈截断把普通微扰误差控制在下一圈阶；若 \(\epsilon_{\log}\gtrsim1\)，固定圈阶中的大对数必须用重整化群或相应因子化结构重求和。Ward--Takahashi 恒等式本身不依赖这一截断。

电子二点函数先给出质量与外腿归一化。把重整化后的约化逆传播子写成
\begin{equation}
{
S_R^{-1}(p)
=\slashed p-m_{\mathrm e} c-\Sigma_R(p).}
\label{eq:qed-renormalized-electron-inverse-dp53}
\end{equation}
一般 Lorentz 协变自能同时含 $\slashed p$ 与标量两种 Dirac 结构，因此记为 $\Sigma_R(p)$。为定义单粒子极点沿质量壳法向的斜率，固定任意未来指向单位类时四矢量 $u^\mu$，满足 $u^2=1$，并沿同一方向写
\begin{equation*}
 p^\mu=\lambda u^\mu,
 \qquad [\lambda]={\rm kg\,m\,s^{-1}}.
\end{equation*}
正能投影算符为
\begin{equation*}
 P_+(u)\equiv\frac{\id(4)+\slashed u}{2},
 \qquad P_+^2=P_+.
\end{equation*}
Lorentz 协变性保证自能在这个二维正能子空间内只留下一个标量本征值；定义
\begin{equation*}
 \sigma_R(\lambda)
 \equiv\frac12\Tr\!\left[P_+(u)\Sigma_R(\lambda u)\right],
 \qquad
 P_+\Sigma_R(\lambda u)P_+=\sigma_R(\lambda)P_+.
\end{equation*}
在壳点 $\lambda=m_{\mathrm e} c$，物理质量要求投影逆传播子在那里有零点，而单位留数要求这个零点对标量极点坐标 $\lambda$ 的一阶斜率等于自由值。两条条件严格写成
\begin{equation}
{
\sigma_R(m_{\mathrm e} c)=0,
\qquad
\left.\frac{\dd\sigma_R(\lambda)}{\dd\lambda}\right|_{\lambda=m_{\mathrm e} c}=0.}
\label{eq:qed-onshell-selfenergy-conditions-dp53}
\end{equation}
第一式固定质量反项，使极点仍在 $p^2=m_{\mathrm e}^2c^2$；第二式固定电子场强反项，使 LSZ 外腿抽出的正能单粒子极点具有单位留数。传统写法
$\Sigma_R(\slashed p)|_{\slashed p=m_{\mathrm e} c}=0$ 与
$\partial\Sigma_R/\partial\slashed p|_{m_{\mathrm e} c}=0$
就是这两个投影标量条件的简写；所谓“对 $\slashed p$ 求导”因此始终指这两个 Lorentz 标量系数的普通导数，而不是未定义的矩阵微分。

光子二点函数则固定电荷的低能归一化。真空极化的横向分解为
\begin{equation*}
\Pi^{\mu\nu}(q)
=(q^2\eta^{\mu\nu}-q^\mu q^\nu)\Pi(q^2),
\end{equation*}
在壳电荷方案中取
\begin{equation}
{\Pi_R(0)=0.}
\label{eq:qed-onshell-vacpol-condition-dp53}
\end{equation}
于是 \(q^2\to0\) 的长程 Coulomb/Thomson 极限直接使用实验测得的低能电荷；非零 \(q^2\) 时剩余的 \(\Pi_R(q^2)\) 才描述真空屏蔽随分辨尺度改变的物理效应。

电磁顶角把前两个问题连接起来。对在壳电子外腿，完整顶角由两个无量纲形状因子参数化：
\begin{equation*}
\bar u(p')\Gamma_R^\mu(p',p)u(p)
=\bar u(p')\left[
F_1(q^2)\gamma^\mu
+\frac{\ii\sigma^{\mu\nu}q_\nu}{2m_{\mathrm e} c}F_2(q^2)
\right]u(p).
\end{equation*}
低能电荷的定义要求
\begin{equation}
{F_1(0)=1,}
\label{eq:qed-F1-charge-normalization-dp53}
\end{equation}
而 \(F_2(0)\) 是 Ward--Takahashi 恒等式没有固定的有限横向结构；它给出反常磁矩
\begin{equation*}
a_e=\frac{g_{\rm mag}-2}{2}=F_2(0).
\end{equation*}
因此质量、低能电荷和反常磁矩分别由电子极点、零动量光子交换和横向顶角结构定义，对应传播子和顶角的三个独立可测信息。

精确 Ward--Takahashi 恒等\eqnrefs{eq:qed_ward_takahashi_momentum} 在一圈阶把电子自能与零动量转移顶角核直接联系起来。两者写为
\begin{equation}
{
\begin{aligned}
\Sigma_{\rm 1L}(p)
&=-\ii g_{\rm em}^2
\int_{\ell}^{(d)}
\frac{\gamma^\alpha\mathsf S_0(p-\ell)\gamma_\alpha}
{\ell^2+\ii0^+},\\
\Lambda_{\rm 1L}^{\mu}(p,p)
&=-\ii g_{\rm em}^2
\int_{\ell}^{(d)}
\frac{\gamma^\alpha\mathsf S_0(p-\ell)\gamma^\mu
\mathsf S_0(p-\ell)\gamma_\alpha}
{\ell^2+\ii0^+},
\end{aligned}}
\label{eq:qed-one-loop-selfenergy-vertex-kernels-dp68}
\end{equation}
其中 \(\mathsf S_0(k)\) 采用自由电子约化传播核及其既定极点处方。两式的区别只有内部电子线上是否插入一个零动量光子顶角。对同一内部传播核关于外动量求导，正好产生这个插入，因此


\begin{proposition}[保持平移的正则化下的 Ward 核恒等式]
设正则化允许同一圈动量平移并保持矢量规范对称，自能与顶角使用同一积分与反项处方。则有
\begin{equation}
{
\Lambda_{\rm 1L}^{\mu}(p,p)
=-\frac{\partial\Sigma_{\rm 1L}(p)}{\partial p_\mu}.}
\label{eq:qed-one-loop-ward-kernel-dp53}
\end{equation}
\end{proposition}
\begin{proof}
正文\eqnrefs{eq:qed-one-loop-selfenergy-vertex-kernels-dp68} 已把零动量转移处的自能图和顶角图写成同一内部电子传播核的两个积分。关键步骤是证明：对自能图的内部电子传播核关于外动量求导，恰好等价于在该电子线上插入一个 \(\gamma^\mu\) 顶角。

自由约化传播核满足
\begin{equation*}
(\slashed k-m_{\mathrm e} c)\mathsf S_0(k)=\id(4).
\end{equation*}
对一般矩阵恒等式 \(M M^{-1}=I\) 求导并左乘 \(M^{-1}\)，得到
\begin{equation*}
\partial_xM^{-1}
=-M^{-1}(\partial_xM)M^{-1}.
\end{equation*}
取 \(M(p)=\slashed p-\slashed\ell-m_{\mathrm e} c\)，则
\begin{equation*}
\frac{\partial\mathsf S_0(p-\ell)}{\partial p_\mu}
=-\mathsf S_0(p-\ell)\gamma^\mu\mathsf S_0(p-\ell).
\end{equation*}
因此对正文自能核逐项求导可得
\begin{align*}
\frac{\partial\Sigma_{\rm 1L}(p)}{\partial p_\mu}
&=-\ii g_{\rm em}^2
\int_{\ell}^{(d)}
\frac{\gamma^\alpha
[\partial\mathsf S_0(p-\ell)/\partial p_\mu]
\gamma_\alpha}
{\ell^2+\ii0^+}\\
&=+\ii g_{\rm em}^2
\int_{\ell}^{(d)}
\frac{\gamma^\alpha\mathsf S_0(p-\ell)
\gamma^\mu\mathsf S_0(p-\ell)\gamma_\alpha}
{\ell^2+\ii0^+}.
\end{align*}
第二行与正文\eqnrefs{eq:qed-one-loop-selfenergy-vertex-kernels-dp68} 中的顶角积分逐因子相同，只差整体负号，所以
\begin{equation*}
\Lambda_{\rm 1L}^{\mu}(p,p)
=-\frac{\partial\Sigma_{\rm 1L}(p)}{\partial p_\mu},
\end{equation*}
即正文\eqnrefs{eq:qed-one-loop-ward-kernel-dp53}。
\end{proof}

这是一圈 Ward--Takahashi 恒等式的图级形式：电子自能对外动量的变化与零动量电磁顶角不是两个独立的紫外结构。

自能中乘在 \(\slashed p\) 前的紫外发散因此必须与最小耦合顶角的纵向紫外发散具有同一系数。若 \(Z_2\) 与 \(Z_1\) 分别表示电子场和顶角的重整化常数，则


\begin{proposition}[Ward 恒等式约束重整化常数]
采用保持矢量规范对称的正则化与统一重整化方案，电子二点与顶角反项按正文定义。则有
\begin{equation}
{Z_1=Z_2.}
\label{eq:qed-Z1-Z2-dp53}
\end{equation}
\end{proposition}
\begin{proof}
\eqnrefs{eq:qed-one-loop-ward-kernel-dp53} 已经说明一圈电子自能的动量依赖与零动量顶角修正并非独立。把同一关系施加到重整化后的二点函数和顶角后，剩余的局域反项也必须满足相同的系数关系。

把电子二点与电磁顶角的反项分别写成
\begin{align*}
S_R^{-1}(p)
&=\slashed p-m_{\mathrm e} c-\Sigma_{\rm 1L}(p)
+\delta_2\slashed p-\delta_m c,\\
\Gamma_R^\mu(p,p)
&=\gamma^\mu+\Lambda_{\rm 1L}^\mu(p,p)
+\delta_1\gamma^\mu.
\end{align*}
对第一式关于 \(p_\mu\) 求导，质量反项不贡献：
\begin{equation*}
\frac{\partial S_R^{-1}}{\partial p_\mu}
=\gamma^\mu
-\frac{\partial\Sigma_{\rm 1L}}{\partial p_\mu}
+\delta_2\gamma^\mu.
\end{equation*}
零动量 Ward 恒等式给出
\begin{equation*}
\Gamma_R^\mu(p,p)
=\frac{\partial S_R^{-1}(p)}{\partial p_\mu}.
\end{equation*}
再用正文\eqnrefs{eq:qed-one-loop-ward-kernel-dp53} 比较两边，树级项和一圈图项已经逐项相同，因此只能有
\begin{equation*}
\delta_1\gamma^\mu=\delta_2\gamma^\mu.
\end{equation*}
于是 \(\delta_1=\delta_2\)，等价地
\begin{equation*}
Z_1=Z_2,
\end{equation*}
得到正文\eqnrefs{eq:qed-Z1-Z2-dp53}。
\end{proof}


电荷关系还需要比较同一个裸相互作用在裸场与重整化场中的两种写法。由
\(\psi_0=Z_2^{1/2}\psi_R\)、\(A_{0\mu}=Z_3^{1/2}A_{R\mu}\)，裸最小耦合项与以 \(Z_1\) 参数化的重整化顶角具有相同系数，当且仅当
\begin{equation}
{
q_0 Z_2 Z_3^{1/2}=q_R Z_1.}
\label{eq:qed-bare-renormalized-charge-match-dp68}
\end{equation}
代入\eqnrefs{eq:qed-Z1-Z2-dp53} 后，电子场因子约去，得到


\begin{proposition}[规范协变的电荷重整化关系]
设裸与重整化场按正文线性缩放，并使用同一 Ward 恒等式与电荷顶角约定。则有
\begin{equation}
{q_0=q_RZ_3^{-1/2}.}
\label{eq:qed-charge-Z3-dp53}
\end{equation}
\end{proposition}
\begin{proof}
关键步骤不是把 \(Z_1=Z_2\) 代进一个已经写好的等式，而是先说明该等式本身从哪里来。裸 QED 的电子--光子相互作用写成
\begin{equation*}
\mathcal L_{\rm int}^{(0)}
=-q_0\,\bar\psi_0\gamma^\mu A_{0\mu}\psi_0.
\end{equation*}
定义重整化场
\begin{equation*}
\psi_0=Z_2^{1/2}\psi_R,
\qquad
\bar\psi_0=Z_2^{1/2}\bar\psi_R,
\qquad
A_{0\mu}=Z_3^{1/2}A_{R\mu},
\end{equation*}
则同一个裸相互作用在重整化场变量中成为
\begin{align*}
\mathcal L_{\rm int}^{(0)}
&=-q_0
\bigl(Z_2^{1/2}\bar\psi_R\bigr)
\gamma^\mu
\bigl(Z_3^{1/2}A_{R\mu}\bigr)
\bigl(Z_2^{1/2}\psi_R\bigr)\\
&=-q_0Z_2Z_3^{1/2}
\bar\psi_R\gamma^\mu A_{R\mu}\psi_R.
\end{align*}
另一方面，重整化拉氏量把同一个局域 Dirac 顶角写成“重整化耦合 + 顶角反项”，即
\begin{equation*}
\mathcal L_{\rm int}^{R+\rm ct}
=-q_R Z_1\,
\bar\psi_R\gamma^\mu A_{R\mu}\psi_R.
\end{equation*}
两式描述的是同一个裸理论，只是变量不同；因为
\(\bar\psi_R\gamma^\mu A_{R\mu}\psi_R\)
是同一个局域算符，其系数必须逐项相等，所以
\begin{equation*}
q_0Z_2Z_3^{1/2}=q_RZ_1,
\end{equation*}
这才得到正文\eqnrefs{eq:qed-bare-renormalized-charge-match-dp68}。

上一推导由 Ward--Takahashi 恒等式得到 \(Z_1=Z_2\)。因此
\begin{equation*}
q_0Z_2Z_3^{1/2}=q_RZ_2.
\end{equation*}
场重整化必须是可逆变量重定义，所以 \(Z_2\neq0\)，约去以后得到
\begin{equation*}
q_0Z_3^{1/2}=q_R,
\qquad
q_0=q_RZ_3^{-1/2},
\end{equation*}
即正文\eqnrefs{eq:qed-charge-Z3-dp53}。

这一步的结构意义现在是可见的：电子外腿的两个 \(Z_2^{1/2}\) 本来确实进入裸顶角，但局域 \(U(1)\) 对称性又强迫顶角反项满足 \(Z_1=Z_2\)，两者严格抵消；因而独立的电荷紫外重整化只能由光子二点函数中的 \(Z_3\) 承担。
\end{proof}

所以阿贝尔电荷重整化的独立紫外信息只剩光子场因子 \(Z_3\)：电子外腿与纵向顶角的紫外重整化已经被局域 \(U(1)\) 对称性锁定为同一个量。

若把一圈 QED 看成对树级理论的第一次形变，可以用三个低能实验量给每个基本核定位。电子静止能确定精确传播子的极点位置 \(E=m_{\mathrm e} c^2\)，因此电子自能的常数部分不能独立成为新参数，而被质量反项重新吸收到物理质量定义中。低动量 Coulomb 或 Thomson 散射确定 \(q_R\)，因此真空极化在 \(q^2=0\) 的常数部分被条件 \(\Pi_R(0)=0\) 吸收；有限 \(Q\) 的偏离才表现为运行耦合。电子自旋在弱磁场中的进动则测量 \(g_{\rm mag}\)，其树级值 \(2\) 由 Dirac 方程给出，而一圈顶角的有限 Pauli 部分给出
\[
F_2(0)=\frac{\alpha}{2\pi}+O(\alpha^2).
\]
这三个观测量分别钉住极点、低能电荷和横向磁矩结构，因此一圈重整化的物理输入与三个一粒子不可约基本核的职责一一对应。
